\documentclass[10pt,a4paper]{amsart}
\usepackage[margin=19mm]{geometry}
\usepackage{amsmath,amssymb,mathtools}
\usepackage[scr=rsfs,cal=euler]{mathalfa}
\usepackage{xfrac}
\usepackage[unicode,hidelinks,bookmarks=false]{hyperref}
\newcommand{\ie}{i.e.\ }
\newcommand{\cf}{cf.\ }
\newcommand{\resp}{respectively\ }
\newcommand{\Subsection}{\subsection}
\providecommand{\nobreakdash}{\nobreak\hbox{-}}
\setlength{\emergencystretch}{2em}
\multlinegap0pt
\usepackage{siunitx}
\usepackage{siunitx}
\usepackage{siunitx}
\usepackage{siunitx}
\usepackage{amssymb}
\usepackage{mathtools}
\usepackage{xcolor}
\usepackage{siunitx}
\usepackage{float}
\usepackage{tikz}
\newtheorem{proposition}{Proposition}
\title{Sustained Limit Cycles in the Logistic Two-Gene Genetic Oscillator: A Delay-Driven Hopf Bifurcation}
\author{Ismail Belgacem}
\date{Source: arXiv:2605.23722v1 (CC BY 4.0)}
\begin{document}
\maketitle

\noindent\emph{Source and licence.} Sustained Limit Cycles in the Logistic Two-Gene Genetic Oscillator: A Delay-Driven Hopf Bifurcation. Version arXiv:2605.23722v1 (CC BY 4.0), distributed by arXiv under the Creative Commons Attribution 4.0 International licence, http://creativecommons.org/licenses/by/4.0/, as recorded in the arXiv licence metadata for this exact version and shown on the abstract page https://arxiv.org/abs/2605.23722v1. Full licence notice: \texttt{licenses/arxiv-2605.23722-CC-BY-4.0.txt}.\par\smallskip

\noindent\textit{Editorial scope.} Counted unit: the article's proposition prop:Cinfty (source0.tex lines 1657--1692). The article's complete original proof of it is delivered with it (source0.tex lines 1693--1744, one \texttt{proof} environment of 52 lines). No mathematics is written, repaired or re-derived: the statement and the proof are the article's own lines, in the article's own notation, and no retained line is substituted or re-flowed.  No further source-local context is needed: every cross-reference of every retained line resolves inside the two delivered spans.  Nothing was omitted from any retained span, so the package carries no omission record.  No retained line contains a citation, so the wrapper carries no bibliography and no bibliography entry.  No retained line contains a figure environment, an image-inclusion command or drawing code, so no image is delivered and the package declares no asset dependency.  Editorial formatting only, in addition: an AMS \texttt{amsart} preamble that restates the article's own declarations and macros used by the retained lines, namely the environments \texttt{proposition} on separate counters, together with the standard packages those lines need.  The retained environments print in this document as Proposition 1 in order of appearance, whereas the article prints the same items with the numbering of its own full document.  The complete native source of the article as distributed in the arXiv e-print for this exact version is bundled unchanged as \texttt{sources/201231/source0.tex}.
\begin{proposition}[Closed-form deep-relaxation period offset]
\label{prop:Cinfty}
Let $M_i:=\kappa_i/\gamma_i$ and assume $0<\theta_i<M_i$ for $i=1,2$.
Assume the deep-relaxation regime: $\lambda\to\infty$ and $\tau\to\infty$
in such a way that, in the periodic regime, each variable reaches a small
neighbourhood of its high or low saturation level (within $O(e^{-c\lambda})$
of $\{0,M_i\}$ for some $c>0$) before the next delayed signal triggers a
switch. Then the period of the limit cycle satisfies
\begin{equation}
T(\tau) \,=\, 2\tau \,+\, C_\infty \,+\, O(e^{-c\lambda})
\qquad(\tau\to\infty),
\label{eq:Cinfty_asymp}
\end{equation}
where the offset $C_\infty$ has the closed form
\begin{equation}
C_\infty
\;=\; \frac{1}{\gamma_1}\,\ln\!\frac{M_1^{2}}{\theta_1(M_1-\theta_1)}
   \,+\,\frac{1}{\gamma_2}\,\ln\!\frac{M_2^{2}}{\theta_2(M_2-\theta_2)}
\;=\; \sum_{i=1}^{2}\frac{1}{\gamma_i}\,
       \ln\!\frac{M_i^{2}}{\theta_i\,(M_i-\theta_i)}.
\label{eq:Cinf}
\end{equation}
The sum decomposes naturally into four threshold-crossing
contributions, $C_\infty = \Delta_A+\Delta_B+\Delta_C+\Delta_D$, where
\begin{equation}
\begin{aligned}
\Delta_A&=\frac{\ln[M_2/(M_2-\theta_2)]}{\gamma_2},
&\quad
\Delta_B&=\frac{\ln[M_1/\theta_1]}{\gamma_1},\\[4pt]
\Delta_C&=\frac{\ln[M_2/\theta_2]}{\gamma_2},
&\quad
\Delta_D&=\frac{\ln[M_1/(M_1-\theta_1)]}{\gamma_1}.
\end{aligned}
\label{eq:phase_durations}
\end{equation}
\end{proposition}

\begin{proof}
In the deep-relaxation regime the limit cycle decomposes into four
phases A, B, C, D distinguished by the sign pattern of the delayed
signals
$(\mathrm{sgn}(x_1(t-\tau_1)-\theta_1),
  \mathrm{sgn}(x_2(t-\tau_2)-\theta_2))$. Within each phase, the
saturated logistic delivers a quasi-constant input ($f^+\approx 1$ or
$0$) to one gene's production, so that gene relaxes exponentially
toward its on or off attractor.

\emph{Phase A (rise of $x_2$).} At the start, $x_1(t-\tau_1)$ has just
crossed $\theta_1$ from below, so $f^+(x_1(t-\tau_1))\approx 1$, which
gives $\dot{x}_2 \approx \kappa_2-\gamma_2 x_2$. Starting from
$x_2\approx 0$ (deep-relaxation hypothesis), the solution is
$x_2(t)\approx M_2\bigl(1-e^{-\gamma_2(t-t_A)}\bigr)$ where $t_A$ is the
entry time. The phase ends when $x_2$ crosses $\theta_2$ from below,
which solves $1-e^{-\gamma_2\Delta_A} = \theta_2/M_2$ for the duration
$\Delta_A=\gamma_2^{-1}\ln[M_2/(M_2-\theta_2)]$.

\emph{Phase B (fall of $x_1$).} After Phase A, $x_2(t-\tau_2)$ crosses
$\theta_2$ from below, switching $f^-(x_2(t-\tau_2))$ from $\approx 1$
to $\approx 0$ and triggering exponential decay
$\dot{x}_1\approx -\gamma_1 x_1$ starting from $x_1\approx M_1$. The
phase ends when $x_1$ crosses $\theta_1$ from above, after
$\Delta_B=\gamma_1^{-1}\ln[M_1/\theta_1]$.

\emph{Phase C (fall of $x_2$).} Symmetric to Phase A, with $x_1$ in the
off state: $x_2$ decays from $\approx M_2$ toward $0$ and crosses
$\theta_2$ from above after $\Delta_C=\gamma_2^{-1}\ln[M_2/\theta_2]$.

\emph{Phase D (rise of $x_1$).} Symmetric to Phase B: $x_1$ rises from
$\approx 0$ toward $M_1$ and crosses $\theta_1$ from below after
$\Delta_D=\gamma_1^{-1}\ln[M_1/(M_1-\theta_1)]$. After Phase D the
configuration returns to the start of Phase A and the cycle repeats.

Each transition between phases is delayed by $\tau_1$ or $\tau_2$
relative to the threshold crossing that triggers it, contributing a
total of $\tau_1+\tau_2+\tau_1+\tau_2 = 2\tau$ to the period.
Summing yields
\begin{equation*}
T(\tau) \,=\, 2\tau + \Delta_A + \Delta_B + \Delta_C + \Delta_D
\,+\, O(e^{-c\lambda}),
\end{equation*}
with the exponential remainder accounting for the
saturation-approach error: the variables enter each phase not exactly
at $0$ or $M_i$ but at a distance $O(e^{-c\lambda})$ from saturation,
where the constant $c$ is determined by the worst-case approach margin
$\min_i\{\theta_i,M_i-\theta_i\}/(M_i)$. Using the elementary identity
$\ln[a/(a-b)]+\ln[a/b]=\ln[a^2/(b(a-b))]$ with $a=M_i$, $b=\theta_i$
to combine $\Delta_A+\Delta_C$ and $\Delta_B+\Delta_D$ yields the
displayed closed form for $C_\infty$.
\end{proof}

\end{document}
